\title{Group Sequence Policy Optimization}

\begin{abstract}
We present Group Sequence Policy Optimization (GSPO), a reinforcement learning algorithm specifically designed to train large language models with greater stability, efficiency, and effectiveness. In contrast to earlier methods that focus on importance ratios at the token level, GSPO employs a sequence-level importance ratio, implementing its clipping, rewarding, and optimization mechanisms across entire sequences. Our empirical results show that GSPO not only outperforms the GRPO algorithm in terms of training speed and result quality, but also brings considerable stability to Mixture-of-Experts (MoE) RL training and offers the prospect of a more streamlined RL infrastructure. The integration of GSPO has been instrumental in advancing the capabilities of the most recent Qwen3 models.
\end{abstract}

\section{Introduction}
\label{sec:intro}

Reinforcement learning (RL) has become a central methodology for advancing the scaling of language models \citep{o1,dpsk-r1,qwq32b,qwen3}. By leveraging large-scale RL, these models acquire enhanced abilities, enabling them to address complex tasks like competitive mathematics and programming via extended and more intricate reasoning pathways.

Achieving successful scaling of RL with increased computational resources hinges primarily on ensuring stable and reliable training dynamics. Nevertheless, existing leading RL algorithms, such as GRPO \citep{grpo}, encounter significant stability challenges when deployed for training extremely large language models; these challenges frequently culminate in catastrophic, unrecoverable collapse of the models \citep{qwen3, minimax-m1}. This pronounced instability presents a substantial obstacle to further expanding the capabilities of language models through additional RL-based training.

In this work, we demonstrate that the instability observed in GRPO arises from an inherent misuse and subsequent breakdown of importance sampling weights within its algorithmic framework. This results in significant training noise of high variance, which escalates as response sequences lengthen and is exacerbated by the use of clipping, thereby leading to eventual model collapse.

To overcome these fundamental shortcomings, we introduce \textbf{Group Sequence Policy Optimization (GSPO)}, a novel reinforcement learning algorithm tailored for large language model training. GSPO’s principal advancement is its theoretically consistent formulation of the importance ratio, which is derived from sequence likelihood \citep{click}, thereby adhering to the foundational principles of importance sampling. Further, GSPO evaluates normalized rewards as advantages calculated across multiple responses to the same query, thus promoting coherence between rewards at the sequence level and the optimization objective.

Through empirical analysis, we find that GSPO notably surpasses GRPO in terms of training stability, efficiency, and overall effectiveness. Importantly, GSPO naturally addresses stability issues present during RL training of large Mixture-of-Experts (MoE) models, obviating the requirement for intricate stabilization approaches and suggesting opportunities for streamlining RL system architecture. The advantages afforded by GSPO have directly led to substantial performance gains in the most recent Qwen3 models. We anticipate that GSPO will serve as a stable and scalable basis for further progress in large-scale RL training involving language models.

\section{Preliminaries}

\paragraph{Notation}
Throughout this work, we refer to an autoregressive language model with parameters $\theta$ as the policy $\pi_\theta$. Here, $x$ represents a query and $\mathcal{D}$ is used for the collection of queries. The probability of generating a response $y$ to a query $x$ via the policy $\pi_\theta$ is expressed as $\pi_\theta(y | x) = \prod_{t=1}^{|y|} \pi_\theta(y_t | x, y_{<t})$, where $|y|$ specifies the total number of tokens in the response $y$. The combination $(x, y)$—representing a query paired with its response—can be evaluated by a verifier function $r$, assigning a reward $r(x, y)$ that falls within the interval $[0, 1]$.

\paragraph{Proximal Policy Optimization (PPO)} Using samples generated from the old policy $\pi_{\theta_\text{old}}$, PPO \citep{ppo} constrains the policy update within a proximal region of the old policy through the clipping mechanism. Specifically, PPO employs the following objective for policy optimization (we omit the KL regularization term hereinafter for brevity, as it is not the focus of this paper):

\begin{align}
\mathcal{J}_\text{PPO}(\theta) = \mathbb{E}_{ x \sim \mathcal{D},\, y \sim \pi_{\theta_\text{old}}( \cdot | x) }
\left[ \frac{1}{|y|} \sum_{t=1}^{|y|} 
\min \left( w_{t}(\theta) \widehat{A}_{t},  \, \mathrm{clip} \left( w_{t}(\theta), 1 - {\varepsilon}, 1 + {\varepsilon}\right) \widehat{A}_{t} \right)
\right],
\end{align}

The token $y_t$'s importance ratio is expressed as $
w_{t}(\theta) = \frac{ \pi_{\theta} (y_{t} | x, y_{<t}) }{ \pi_{\theta_\text{old}} (y_{t} | x,y_{<t})} $, with the advantage $\widehat{A}_{t}$ for $y_t$ being provided by a separate value model, and $\varepsilon$ serving as the threshold for importance ratio clipping.

A fundamental obstacle encountered by PPO in practical applications is its significant dependence on the value model. Typically, the value model matches the policy model in scale, thereby imposing substantial demands on both memory capacity and computation. Moreover, the effectiveness of PPO is predicated on the quality of the value model's estimates. Developing a robust value model is difficult in itself, but achieving scalability for longer output sequences and handling more intricate tasks is even more problematic.

\paragraph{Group Relative Policy Optimization (GRPO)} GRPO \citep{grpo} bypasses the need for the value model by computing the relative advantage of each response within a group of responses to the same query. Specifically, GRPO optimizes the following objective:

\begin{align}
\label{equ:grpo}
\mathcal{J}_\text{GRPO}(\theta) = \mathbb{E}_{ x \sim \mathcal{D},\, \{y_i\}_{i=1}^G \sim \pi_{\theta_\text{old}}( \cdot | x) }
\left[ \frac{1}{G} \sum_{i=1}^{G} \frac{1}{|y_i|} \sum_{t=1}^{|y_i|} 
\min \left( w_{i,t}(\theta) \widehat{A}_{i,t},  \, \mathrm{clip} \left( w_{i,t}(\theta), 1 - {\varepsilon}, 1 + {\varepsilon}\right) \widehat{A}_{i,t} \right)
\right],
\end{align}

where $G$ is the number of generated responses to each query $x$ (i.e., the group size), and the importance ratio $w_{i,t}(\theta)$ and advantage $\widehat{A}_{i,t}$ of token $y_{i,t}$ are:

\begin{align}
    w_{i,t}(\theta)=\frac{ \pi_{\theta} (y_{i,t} | x, y_{i,<t}) }{ \pi_{\theta_\text{old}} (y_{i,t} | x,y_{i,<t})},\quad\ 
    \widehat{A}_{i,t} = \widehat{A}_{i} = \frac{r(x, y_i) - \mathrm{mean} \left( \{ r(x, y_i) \}_{i=1}^G \right) }{ \mathrm{std} \left( \{ r(x, y_i) \}_{i=1}^G \right) },
\end{align}

respectively, where all the tokens in $y_i$ share the same advantage as $\widehat{A}_{i}$.

\section{Motivation}
\label{sec:motivation}

The increasing scale of models, higher sparsity as observed in Mixture-of-Experts architectures, and the tendency toward longer generated responses require employing large rollout batch sizes to achieve efficient hardware utilization during reinforcement learning. To enhance sample efficiency, practitioners typically split this extensive batch of rollout trajectories into several mini-batches, which are then used for computing gradient updates. This batching approach inherently results in an off-policy learning context: the responses $y$ are drawn from a prior policy $\pi_{\theta_\text{old}}$ rather than the updated policy $\pi_\theta$ currently being trained. This aspect underpins the need for clipping mechanisms, such as those implemented in PPO and GRPO, which restrict the contribution of samples that are excessively ``off-policy'' during gradient estimation.

While mechanisms like clipping aim to manage this off-policy discrepancy, we identify a more fundamental issue in GRPO: \textit{its objective is ill-posed}. This problem becomes particularly acute when training large models on long-response tasks, leading to catastrophic model collapse. The ill-posed nature of the GRPO objective stems from a misapplication of importance sampling weights. The principle of importance sampling is to estimate the expectation of a function $f$ under a target distribution $\pi_\text{tar}$ by re-weighting samples drawn from a behavior distribution $\pi_\text{beh}$:

\begin{align}
\mathbb{E}_{z \sim \pi_\text{tar}} \left[ f(z) \right] = \mathbb{E}_{z \sim \pi_\text{beh}} \left[ \frac{ \pi_\text{tar} (z) }{ \pi_\text{beh} (z)} \, f(z) \right].
\end{align}

This approach fundamentally depends on aggregating over a large number of samples ($N \gg 1$) drawn from the behavior distribution $\pi_\text{beh}$, allowing the importance weight $\frac{ \pi_\text{tar} (z) }{ \pi_\text{beh} (z)}$ to properly mitigate the distributional discrepancy.

By contrast, GRPO incorporates the importance weight $\frac{ \pi_{\theta} (y_{i,t} | x, y_{i,<t}) }{ \pi_{\theta_\text{old}} (y_{i,t} | x,y_{i,<t})}$ at every token location $t$. This weighting utilizes a lone sample $y_{i,t}$ from the relevant next-token distribution $\pi_{\theta_\text{old}}( \cdot | x, y_{i, <t})$, and thus cannot execute effective distribution correction. Rather, it infuses substantial variance into the gradient estimates during training; this effect is cumulative for extended sequences and becomes more pronounced due to the clipping procedure. Empirical evidence indicates that this process can precipitate model collapse—a state from which recovery is rarely possible. Once collapsed, attempts to restart training fail even with careful checkpoint restoration, elaborate hyperparameter adjustment (such as varying clipping thresholds), increased generation lengths, or alternative RL querying.

These findings highlight an intrinsic flaw in GRPO’s architecture. The ineffectiveness of token-wise importance weighting signals a fundamental tenet: \textbf{the scale of the optimization objective must be consistent with the scale at which reward is assigned}. Since reward operates at the sequence level, token-based off-policy adjustment is problematic. Accordingly, we propose abandoning token-level formulations in favor of leveraging importance weighting and conducting optimization at the \textit{sequence level}.

\section{Algorithm}

\subsection{GSPO: Group Sequence Policy Optimization}

Although GRPO faces issues with the token-level importance weight $\frac{ \pi_{\theta} (y_{i,t} | x, y_{i,<t}) }{ \pi_{\theta_\text{old}} (y_{i,t} | x,y_{i,<t})}$, we find that the \textit{sequence-level} importance weight $\frac{ \pi_{\theta} (y | x) }{ \pi_{\theta_\text{old}} (y | x)}$ in language generation carries a well-defined theoretical interpretation. Specifically, it measures the degree to which a sequence $y$ generated from $\pi_{\theta_\text{old}} (\cdot | x)$ diverges from the updated distribution $\pi_{\theta} (\cdot | x)$, effectively corresponding to the sequence-level reward and functioning as a meaningful proxy for the clipping effect.

Based on this straightforward observation, we propose the \textbf{Group Sequence Policy Optimization (GSPO)} algorithm. GSPO employs the following sequence-level optimization objective:

\begin{align}
\mathcal{J}_\text{GSPO} (\theta) =
\mathbb{E}_{ x \sim \mathcal{D},\, \{y_i\}_{i=1}^G \sim \pi_{\theta_\text{old}}( \cdot | x) }
\left[ 
\frac{1}{G} \sum_{i=1}^{G}
\min \left( s_{i}(\theta)  \widehat{A}_{i},  \, \mathrm{clip} \left( s_{i}(\theta), 1 - {\varepsilon}, 1 + {\varepsilon} \right) \widehat{A}_{i} \right) 
\right],
\label{equ:gspo}
\end{align}

where we adopt the group-based advantage estimation:

\begin{align}
\widehat{A}_{i} = \frac{r(x, y_i) - \mathrm{mean} \left( \{ r(x, y_i) \}_{i=1}^G \right) }{ \mathrm{std} \left( \{ r(x, y_i) \}_{i=1}^G \right) },
\end{align}

and define the importance ratio $s_{i}(\theta)$ based on sequence likelihood \citep{click}:

\begin{align}
s_{i}(\theta) = \left( \frac{ \pi_{\theta} (y_i | x) }{ \pi_{\theta_\text{old}} (y_i | x)} \right)^{\frac{1}{|y_i|}}
=
\exp \left( \frac{1}{|y_i|} \sum_{t=1}^{|y_i|} \log \frac{ \pi_{\theta} (y_{i,t} | x, y_{i,<t}) }{ \pi_{\theta_\text{old}} (y_{i,t} | x,y_{i,<t})} \right).
\end{align}

As a result, GSPO performs clipping on whole responses, rather than on individual tokens, so as to filter out samples that are excessively “off-policy” from the gradient estimation process. This approach aligns with both the reward assignment at the sequence level and the overall optimization. Furthermore, we employ length normalization in $s_{i}(\theta)$, which serves to mitigate variance and ensures that $s_{i}(\theta)$ stays within a consistent numerical interval. Without such normalization, alterations in token likelihoods for just a few tokens could lead to significant variability in the sequence-level importance ratio, necessitating different clipping boundaries for responses of varying lengths. Additionally, it is important to highlight that the clipping intervals used in GSPO, as compared to prior algorithms such as GRPO, often vary substantially in magnitude, a consequence of the differing formulations of importance ratios.

\subsection{Gradient Analysis}
\label{subsec:gradient}

We can derive the gradient of the GSPO objective as follows (clipping is omitted for brevity):

\begin{align}
\nabla_{\theta} \mathcal{J}_\text{GSPO} (\theta)
=&\ 
\nabla_{\theta} \mathbb{E}_{ x \sim \mathcal{D},\, \{y_i\}_{i=1}^G \sim \pi_{\theta_\text{old}}( \cdot | x) }
\left[ 
\frac{1}{G} \sum_{i=1}^{G} s_{i}(\theta) \widehat{A}_{i}
\right] \\
= &\ 
\mathbb{E}_{ x \sim \mathcal{D},\, \{y_i\}_{i=1}^G \sim \pi_{\theta_\text{old}}( \cdot | x) }
\left[ 
\frac{1}{G} \sum_{i=1}^{G}
s_{i}(\theta) \widehat{A}_{i}
\cdot \nabla_{\theta} \log s_{i}(\theta)
\right] \\
=&\ 
\mathbb{E}_{ x \sim \mathcal{D},\, \{y_i\}_{i=1}^G \sim \pi_{\theta_\text{old}}( \cdot | x) }
\left[ 
\frac{1}{G} \sum_{i=1}^{G} 
\left( \frac{ \pi_{\theta} (y_i | x) }{ \pi_{\theta_\text{old}} (y_i | x)} \right)^{\frac{1}{|y_i|}} \widehat{A}_{i} 
\cdot \frac{1}{|y_i|} \sum_{t=1}^{|y_i|} \nabla_{\theta} \log \pi_{\theta} (y_{i,t} | x, y_{i,<t}) 
\right].
\label{equ:gspo_grad}
\end{align}

For comparison, the gradient of the GRPO objective is as follows (note that $\widehat{A}_{i,t} = \widehat{A}_{i}$):

\begin{align}
\nabla_{\theta} \mathcal{J}_\text{GRPO} (\theta) 
=&\ 
\nabla_{\theta} \mathbb{E}_{ x \sim \mathcal{D},\, \{y_i\}_{i=1}^G \sim \pi_{\theta_\text{old}}( \cdot | x) }
\left[ \frac{1}{G} \sum_{i=1}^{G} \frac{1}{|y_i|} \sum_{t=1}^{|y_i|} 
w_{i,t}(\theta) \widehat{A}_{i,t}
\right] \\
=&\
\mathbb{E}_{ x \sim \mathcal{D},\, \{y_i\}_{i=1}^G \sim \pi_{\theta_\text{old}}( \cdot | x) }
\left[ \frac{1}{G} \sum_{i=1}^{G} \widehat{A}_{i} 
\cdot \frac{1}{|y_i|} \sum_{t=1}^{|y_i|} 
\frac{ \pi_{\theta} (y_{i,t} | x, y_{i,<t}) }{ \pi_{\theta_\text{old}} (y_{i,t} | x,y_{i,<t})} 
\nabla_{\theta} \log \pi_{\theta} (y_{i,t} | x, y_{i,<t})  
\right].
\end{align}

As such, the primary difference between GSPO and GRPO concerns the \textit{method by which each scheme assigns weights to the gradients of the log likelihoods for tokens}. GRPO employs an ``importance weight'' for each token, given by $\frac{ \pi_{\theta} (y_{i,t} | x, y_{i,<t}) }{ \pi_{\theta_\text{old}} (y_{i,t} | x,y_{i,<t})}$, to scale their respective gradients. Notably, these varying weights can assume values within $(0, 1+\varepsilon]$ when $\widehat{A}_{i} >0$, or within $[1-\varepsilon, +\infty)$ when $\widehat{A}_{i} < 0$. Such variability in token weighting is substantial and can accumulate over the course of optimization, potentially introducing instability and unforeseen behaviors. By contrast, GSPO applies uniform weighting to all tokens in a sequence, thereby removing this potential source of instability associated with the gradient updates in GRPO.

\subsection{GSPO-token: A Token-level Objective Variant}

In scenarios like multi-turn RL, we may desire a finer-grained advantage adjustment than the sequence level. To this end, we introduce a token-level objective variant of GSPO, namely \textbf{GSPO-token}, to allow token-wise advantage customization:

\begin{align}
\mathcal{J}_\text{GSPO-token}(\theta)
=
\mathbb{E}_{ x \sim \mathcal{D},\, \{y_i\}_{i=1}^G \sim \pi_{\theta_\text{old}}( \cdot | x) }
\left[ 
\frac{1}{G} \sum_{i=1}^{G} \frac{1}{|y_i|} \sum_{t=1}^{|y_i|} 
\min \left( s_{i,t}(\theta) \widehat{A}_{i,t},  \, \mathrm{clip} \left( s_{i,t}(\theta), 1 - {\varepsilon}, 1 + {\varepsilon} \right) \widehat{A}_{i,t} \right) 
\right],
\label{equ:gspo-token}
\end{align}

where

\begin{align}
s_{i,t}(\theta) 
= \mathrm{sg} \left[ s_{i}(\theta) \right]  \cdot \frac{ \pi_{\theta} (y_{i,t} | x, y_{i,<t}) }{ \mathrm{sg} \left[ \pi_{\theta} (y_{i,t} | x, y_{i,<t}) \right] },
\end{align}

and $\mathrm{sg}[\cdot]$ denotes only taking the numerical value but stopping the gradient, corresponding to the \texttt{detach} operation in PyTorch. The gradient of GSPO-token can be derived as:

\begin{align}
\nabla_{\theta} \mathcal{J}_\text{GSPO-token}(\theta)
=&\ 
\nabla_{\theta} \mathbb{E}_{ x \sim \mathcal{D},\, \{y_i\}_{i=1}^G \sim \pi_{\theta_\text{old}}( \cdot | x) }
\left[ 
\frac{1}{G} \sum_{i=1}^{G}
\frac{1}{|y_i|} \sum_{t=1}^{|y_i|} 
s_{i,t}(\theta) \widehat{A}_{i,t}
\right] \\
=&\ 
\mathbb{E}_{ x \sim \mathcal{D},\, \{y_i\}_{i=1}^G \sim \pi_{\theta_\text{old}}( \cdot | x) }
\left[ 
\frac{1}{G} \sum_{i=1}^{G} s_i (\theta) \cdot \frac{1}{|y_i|} \sum_{t=1}^{|y_i|} 
\widehat{A}_{i,t} \frac{ \nabla_{\theta} \pi_{\theta} (y_{i,t} | x, y_{i,<t}) }{ \pi_{\theta} (y_{i,t} | x, y_{i,<t}) }
\right] \\
=&\ 
\mathbb{E}_{ x \sim \mathcal{D},\, \{y_i\}_{i=1}^G \sim \pi_{\theta_\text{old}}( \cdot | x) }
\left[ 
\frac{1}{G} \sum_{i=1}^{G}  \left( \frac{ \pi_{\theta} (y_i | x) }{ \pi_{\theta_\text{old}} (y_i | x)} \right)^{\frac{1}{|y_i|}}
\cdot \frac{1}{|y_i|} \sum_{t=1}^{|y_i|}
\widehat{A}_{i,t} \nabla_{\theta} \log \pi_{\theta} (y_{i,t} | x, y_{i,<t})  
\right].
\label{equ:gspo-token_grad}
\end{align}

It is important to observe that $\frac{ \pi_{\theta} (y_{i,t} | x, y_{i,<t}) }{ \mathrm{sg} \left[ \pi_{\theta} (y_{i,t} | x, y_{i,<t}) \right] }$ evaluates to 1 for any choice, resulting in $s_{i,t}(\theta)$ being numerically equivalent to $s_{i}(\theta)$. Analyzing Equation~\eqref{equ:gspo} alongside~\eqref{equ:gspo-token}, as well as Equation~\eqref{equ:gspo_grad} compared to~\eqref{equ:gspo-token_grad}, reveals that GSPO-token and GSPO yield identical results for the optimization objective, clipping constraints, and theoretical gradient calculations, provided that all tokens within a response $y_i$ share an identical advantage (specifically, when $\widehat{A}_{i,t} = \widehat{A}_{i}$). Notably, GSPO-token offers additional versatility through the capacity to specify distinct advantages for each individual token.

\section{Experiments and Discussion}

\subsection{Empirical Results}

We utilize a cold-start model initialized from Qwen3-30B-A3B-Base and evaluate both the reward curves obtained during training and the model’s performance on the AIME'24 benchmark (mean Pass@1 with 32 samplings), LiveCodeBench (202410-202502, mean Pass@1 with 8 samplings), and CodeForces (Elo Rating). Throughout RL training, we divide each batch of rollout data into four separate mini-batches for performing gradient updates. For GSPO, the left and right clipping intervals in Equation~\eqref{equ:gspo} are set to 3e-4 and 4e-4, respectively. In our baseline comparison, we employ GRPO, where the clipping bounds in Equation~\eqref{equ:grpo} are chosen as 0.2 (left) and 0.27 (right), after thorough tuning to ensure equitable evaluation. Importantly, GRPO depends on the Routing Replay training scheme to achieve proper convergence within MoE RL settings—a topic elaborated in \S~\ref{subsec:routing_replay}—while \textit{GSPO eliminates this requirement}.

As illustrated in Figure~\ref{fig:results}, GSPO enables consistently stable training dynamics. Notably, \textbf{GSPO supports ongoing gains in performance as the training compute increases, the query set is periodically refreshed, and generation lengths are extended}. Additionally, GSPO outperforms GRPO with respect to training efficiency, attaining superior accuracy and benchmark scores under equivalent training compute and with the same number of queries consumed. Lastly, integrating GSPO into RL training for the most recent Qwen3 models has yielded strong evidence of GSPO’s effectiveness at harnessing RL scaling for enhancing large language model capabilities.

\subsection{Curious Observation on Clipping Fractions}

A principal difference between GSPO and GRPO lies in GSPO’s strategy of clipping responses at the sequence level, as opposed to GRPO’s token-wise clipping. Notably, as illustrated in Figure~\ref{fig:clipping}, there exists a two orders of magnitude discrepancy in the proportion of tokens clipped by GSPO versus GRPO (even when the clipping intervals are modified, the difference in scale persists). Nonetheless, even though GSPO clips a far larger percentage of tokens—thus reducing the number of tokens utilized for training or gradient estimation—it consistently yields greater training efficiency than GRPO. This seemingly paradoxical outcome—where extensively clipping tokens still enhances training efficiency—suggests that GRPO’s token-based gradient estimates are subject to substantial noise and are suboptimal for leveraging sampled data. Conversely, the sequence-centric method adopted by GSPO supplies a more stable and potent signal for learning.

\subsection{Benefit of GSPO for MoE Training}
\label{subsec:routing_replay}

\paragraph{Background}
In contrast to the reinforcement learning (RL) training of dense architectures, MoE models—owing to their sparse activation patterns—face distinct challenges regarding training stability. Our experiments indicate that when employing the GRPO algorithm, RL training often fails to converge reliably due to the MoE models' tendency for \textit{expert-activation volatility}. Specifically, the set of experts selected to process the same sample may shift dramatically after one or more gradient updates. For instance, using the 48-layer Qwen3-30B-A3B-Base architecture, each RL gradient update leads to a scenario where approximately 10\% of the experts chosen for a given rollout under the updated policy $\pi_{\theta}$ differ from those selected under the previous policy $\pi_{\theta_\text{old}}$. This instability is exacerbated in deeper MoE configurations, resulting in severe fluctuations in the token-level importance weights $w_{i,t}(\theta) = \frac{ \pi_{\theta} (y_{i,t} | x, y_{i,<t}) }{ \pi_{\theta_\text{old}} (y_{i,t} | x,y_{i,<t})}$. As elaborated in \S~\ref{sec:motivation} and~\ref{subsec:gradient}, such volatility can undermine the validity of importance sampling and ultimately disrupt the proper convergence of RL training.

\paragraph{Our Previous Approach} In prior work, we introduced the \textbf{Routing Replay} method to address this issue. Our approach involves recording the expert routing decisions produced by $\pi_{\theta_\text{old}}$ and subsequently applying these recorded routes in $\pi_{\theta}$ during the calculation of importance weights, specifically $w_{i,t}(\theta) = \frac{ \pi_{\theta} (y_{i,t} | x, y_{i,<t}) }{ \pi_{\theta_\text{old}} (y_{i,t} | x,y_{i,<t})}$. By doing so, both $\pi_{\theta} (y_{i,t} | x, y_{i,<t})$ and $\pi_{\theta_\text{old}} (y_{i,t} | x,y_{i,<t})$ process each token $y_{i,t}$ using an identical activated set of experts. This alignment stabilizes the token-level importance weights and ensures that optimization targets a consistent expert network across gradient steps.
As illustrated in Figure~\ref{fig:routing_replay}, Routing Replay is crucial for achieving stable convergence during GRPO training in MoE architectures.

\paragraph{Benefit of GSPO} While Routing Replay permits GRPO training on MoE models to achieve proper convergence, its reliance on recycled routing patterns results in increased memory and communication demands, and may restrict the MoE model’s effective capacity. By contrast, as illustrated in Figure~\ref{fig:results}, GSPO removes the requirement for Routing Replay. It computes importance ratios $s_i(\theta)$ in the usual manner, facilitating typical convergence and consistent optimization. The central observation here is that GSPO is concerned with sequence-level likelihood (i.e., $\pi_{\theta} (y_i | x)$), rather than token-level likelihoods (i.e., $\pi_{\theta} (y_{i,t} | x, y_{i,<t})$). Because the MoE model consistently retains its language modeling ability, the sequence likelihood remains relatively stable. To conclude, GSPO directly addresses the instability in expert activation present in MoE models, eliminating the necessity for sophisticated strategies such as Routing Replay. This approach streamlines and fortifies training, granting the model access to its complete capacity without imposed limitations.

\subsection{Benefit of GSPO for RL Infrastructure}

Due to the differences in numerical precision between engines used for training, such as Megatron, and those deployed for inference, including SGLang and vLLM, it is standard practice to recalculate the likelihoods of responses sampled by the previous policy $\pi_{\theta_\text{old}}$ using the training engine. In contrast, GSPO relies on sequence-level likelihoods instead of those at the token level during optimization, and sequence-level metrics are, intuitively, less sensitive to such precision mismatches. As a result, GSPO enables practitioners to directly utilize the likelihoods obtained from the inference engine without necessitating recomputation on the training engine. This characteristic is particularly advantageous in contexts such as partial rollouts, multi-turn reinforcement learning, and training-inference separation frameworks.

\section{Conclusion}

In this work, we introduce Group Sequence Policy Optimization (GSPO), a novel reinforcement learning approach tailored for optimizing large language models. GSPO leverages importance sampling at its core, establishing importance ratios grounded in sequence likelihood and integrating sequence-level mechanisms for clipping, reward assignment, and optimization. Empirically, GSPO attains significantly enhanced stability, efficiency, and overall performance relative to GRPO, and it proves especially advantageous for large-scale reinforcement learning applied to MoE architectures. This capability underpins the remarkable advancements observed in the latest Qwen3 models. Serving as a scalable foundation, GSPO enables further scaling of reinforcement learning, and we anticipate that this will facilitate meaningful progress towards more advanced forms of intelligence.

\end{document}